
\chapter{Intersection Forms}
\label{ch:viii34-intersection-forms}

Poincaré duality becomes especially rigid in the middle dimension.  On a
closed oriented manifold, cup product followed by evaluation on the
fundamental class produces a bilinear pairing whose algebra records geometric
intersection.  In dimension four this pairing is one of the basic topological
invariants.

\section{The middle-dimensional pairing}
Let $M$ be a closed oriented $2n$-manifold.  Define
\[
Q_M(\alpha,\beta)
=
\langle \alpha\smile\beta,[M]\rangle,
\qquad
\alpha,\beta\in H^n(M;\mathbb Z)/\mathrm{tors}.
\]

\begin{definition}[Intersection form]
\label{def:viii34-intersection-form}
The integral middle-dimensional intersection form is the bilinear form
\[
Q_M:
\bigl(H^n(M;\mathbb Z)/\mathrm{tors}\bigr)^2\longrightarrow\mathbb Z
\]
given by cup product and evaluation on the fundamental class.
\end{definition}

\section{Symmetry and graded symmetry}
Graded commutativity gives
\[
Q_M(\beta,\alpha)=(-1)^nQ_M(\alpha,\beta).
\]
Thus the middle pairing is symmetric when $n$ is even and skew-symmetric when
$n$ is odd.

\section{Unimodularity}
\begin{theorem}[Unimodularity]
\label{thm:viii34-unimodular}
For a closed oriented manifold, the integral intersection form on the free
middle cohomology is unimodular.
\end{theorem}
This is the lattice form of integral Poincaré duality.

\section{Geometric intersections}
If $\alpha$ and $\beta$ are Poincaré dual to complementary transverse cycles,
then
\[
Q_M(\alpha,\beta)
\]
is their signed intersection number.

\section{Self-intersection}
A middle-dimensional submanifold can be pushed slightly in a normal direction.
Its signed intersection with the push-off is its self-intersection number.
For an oriented normal bundle, this number is the Euler number of the normal
bundle.

\section{Four-manifolds}
For a closed oriented $4$-manifold,
\[
Q_M:H^2(M;\mathbb Z)_{\mathrm{free}}\times
H^2(M;\mathbb Z)_{\mathrm{free}}\to\mathbb Z
\]
is symmetric.

\section{Positive and negative directions}
After tensoring with $\mathbb R$, a symmetric form decomposes into positive
and negative directions.  Write
\[
b_2^+(M),\qquad b_2^-(M)
\]
for their dimensions.

\begin{definition}[Signature]
\label{def:viii34-signature}
\[
\sigma(M)=b_2^+(M)-b_2^-(M).
\]
\end{definition}

\section{Definite and indefinite forms}
The form is positive definite if every nonzero vector has positive square,
negative definite if every nonzero vector has negative square, and indefinite
if both signs occur.

\section{Even and odd forms}
An integral symmetric form is even when
\[
Q(x,x)\in2\mathbb Z
\]
for every integral vector $x$; otherwise it is odd.

\section{The four-sphere}
Since
\[
H^2(S^4;\mathbb Z)=0,
\]
the intersection form of $S^4$ has rank zero.

\section{Complex projective plane}
Let $h\in H^2(\mathbb{CP}^2;\mathbb Z)$ be the standard generator with
\[
\langle h^2,[\mathbb{CP}^2]\rangle=1.
\]
Then
\[
Q_{\mathbb{CP}^2}=[1].
\]
It is positive definite, odd, unimodular, and has signature $1$.

\section{Orientation reversal}
Changing the orientation changes the fundamental class by a sign:
\[
Q_{-M}=-Q_M.
\]
Thus
\[
Q_{-\mathbb{CP}^2}=[-1].
\]

\section{The product $S^2\times S^2$}
Let $a,b\in H^2(S^2\times S^2;\mathbb Z)$ be pulled back from the two factors.
Then
\[
a^2=b^2=0,\qquad
\langle a\smile b,[S^2\times S^2]\rangle=1.
\]
Hence
\[
Q_{S^2\times S^2}
=
H
=
\begin{pmatrix}
0&1\\
1&0
\end{pmatrix}.
\]
The hyperbolic form $H$ is even, indefinite, unimodular, and has signature
zero.

\section{Connected sums}
For oriented closed $4$-manifolds,
\[
Q_{M\#N}\cong Q_M\oplus Q_N.
\]
Consequently
\[
Q_{\mathbb{CP}^2\#(-\mathbb{CP}^2)}
=
\begin{pmatrix}
1&0\\
0&-1
\end{pmatrix}.
\]

\section{Integral change of basis}
If $A$ represents $Q_M$ and $P\in GL_r(\mathbb Z)$ changes the integral basis,
then
\[
A\longmapsto P^TAP.
\]
Thus rank, unimodularity, parity, definiteness, and signature are invariants of
the integral congruence class.

\section{Intersection matrices from geometry}
Choosing geometric representatives for a basis of middle homology turns the
form into an intersection matrix.  Diagonal entries are self-intersections;
off-diagonal entries count pairwise signed intersections.

\section{Bridge to Lefschetz theory}
For a self-map $f:M\to M$, fixed points are intersections of the graph
$\Gamma_f\subset M\times M$ with the diagonal $\Delta\subset M\times M$.
Lefschetz theory converts that global intersection number into a trace on
homology.

\section{Solved dossiers}
\begin{problem}
Define the middle-dimensional intersection form of a closed oriented $2n$-manifold.
\end{problem}
\begin{solution}
On the free part, $Q_M(\alpha,\beta)=\langle \alpha\smile\beta,[M]\rangle$ for $\alpha,\beta\in H^n(M;\mathbb Z)/\mathrm{tors}$.
\end{solution}

\begin{problem}
How does graded commutativity control the symmetry of $Q_M$?
\end{problem}
\begin{solution}
$Q_M(\beta,\alpha)=(-1)^nQ_M(\alpha,\beta)$; it is symmetric for even $n$ and skew-symmetric for odd $n$.
\end{solution}

\begin{problem}
Why is the integral form unimodular on the free part?
\end{problem}
\begin{solution}
Integral Poincaré duality identifies the free middle cohomology lattice with its integral dual, so the representing matrix has determinant $\pm1$.
\end{solution}

\begin{problem}
What is the intersection form of $S^4$?
\end{problem}
\begin{solution}
$H^2(S^4)=0$, so the form has rank zero.
\end{solution}

\begin{problem}
Compute the form of $\mathbb{CP}^2$ with its complex orientation.
\end{problem}
\begin{solution}
If $h\in H^2(\mathbb{CP}^2)$ and $\langle h^2,[\mathbb{CP}^2]\rangle=1$, then the matrix is $[1]$.
\end{solution}

\begin{problem}
What does orientation reversal do to the form?
\end{problem}
\begin{solution}
$[-M]=-[M]$, hence $Q_{-M}=-Q_M$.
\end{solution}

\begin{problem}
Compute the form of $S^2\times S^2$.
\end{problem}
\begin{solution}
With $a,b$ dual to the two factors, $a^2=b^2=0$ and $ab$ evaluates to $1$, so the matrix is $\begin{pmatrix}0&1\\1&0\end{pmatrix}$.
\end{solution}

\begin{problem}
What is the signature of the hyperbolic form?
\end{problem}
\begin{solution}
Its eigenvalues are $1$ and $-1$, so the signature is $0$.
\end{solution}

\begin{problem}
Define positive and negative indices in dimension four.
\end{problem}
\begin{solution}
After extending $Q_M$ to $\mathbb R$, $b_2^+$ and $b_2^-$ are the dimensions of maximal positive- and negative-definite subspaces.
\end{solution}

\begin{problem}
Define the signature.
\end{problem}
\begin{solution}
$\sigma(M)=b_2^+-b_2^-$.
\end{solution}

\begin{problem}
How does connected sum affect the form for oriented closed $4$-manifolds?
\end{problem}
\begin{solution}
On the middle cohomology free parts, $Q_{M\#N}\cong Q_M\oplus Q_N$.
\end{solution}

\begin{problem}
Compute $Q_{\mathbb{CP}^2\#(-\mathbb{CP}^2)}$.
\end{problem}
\begin{solution}
It is $\operatorname{diag}(1,-1)$, hence indefinite of signature zero.
\end{solution}

\begin{problem}
What does even mean for an integral symmetric form?
\end{problem}
\begin{solution}
$Q(x,x)$ is even for every integral lattice vector $x$.
\end{solution}

\begin{problem}
Is the $\mathbb{CP}^2$ form even or odd?
\end{problem}
\begin{solution}
Odd, because the generator has self-intersection $1$.
\end{solution}

\begin{problem}
Is the hyperbolic form even?
\end{problem}
\begin{solution}
Yes: for $(x,y)$ its square is $2xy$, always even.
\end{solution}

\begin{problem}
Give the geometric meaning of $Q_M(\alpha,\beta)$ in dimension four.
\end{problem}
\begin{solution}
Represent the dual classes by oriented embedded or immersed surfaces in general position; the form is their signed intersection number.
\end{solution}

\begin{problem}
What is self-intersection geometrically?
\end{problem}
\begin{solution}
It is the intersection of a surface with a transverse push-off, equivalently the Euler number of its oriented normal bundle.
\end{solution}

\begin{problem}
Why does torsion not enter the ordinary integral matrix directly?
\end{problem}
\begin{solution}
The cup-evaluation pairing to $\mathbb Z$ vanishes on torsion against the free fundamental class, so the lattice form is defined on the free quotient.
\end{solution}

\begin{problem}
How does a change of integral basis affect the matrix?
\end{problem}
\begin{solution}
$A$ changes by integral congruence $A\mapsto P^{T}AP$ with $P\in GL_r(\mathbb Z)$.
\end{solution}

\begin{problem}
Why is the chapter a direct continuation of Poincaré duality?
\end{problem}
\begin{solution}
The intersection form is precisely the middle-degree Poincaré pairing promoted from a duality statement to a structured bilinear invariant.
\end{solution}

\section{Exercises}
\begin{exercise}
Compute the intersection form of $S^4$.
\end{exercise}
\begin{exercise}
Compute the form of $\mathbb{CP}^2$.
\end{exercise}
\begin{exercise}
Compute the form of $-\mathbb{CP}^2$.
\end{exercise}
\begin{exercise}
Compute the form of $S^2\times S^2$.
\end{exercise}
\begin{exercise}
Diagonalize the hyperbolic form over $\mathbb R$.
\end{exercise}
\begin{exercise}
Compute its signature.
\end{exercise}
\begin{exercise}
Show that the hyperbolic form is even.
\end{exercise}
\begin{exercise}
Show that $[1]$ is odd.
\end{exercise}
\begin{exercise}
Compute the form of $\mathbb{CP}^2\#\mathbb{CP}^2$.
\end{exercise}
\begin{exercise}
Compute the form of $\mathbb{CP}^2\#(-\mathbb{CP}^2)$.
\end{exercise}
\begin{exercise}
Determine $b_2^+$ and $b_2^-$ in the preceding two examples.
\end{exercise}
\begin{exercise}
Verify that orientation reversal negates the matrix.
\end{exercise}
\begin{exercise}
Show that integral change of basis acts by congruence.
\end{exercise}
\begin{exercise}
Show that unimodularity is basis independent.
\end{exercise}
\begin{exercise}
Derive symmetry in dimension four from graded commutativity.
\end{exercise}
\begin{exercise}
Derive skew-symmetry in middle degree on a closed oriented surface.
\end{exercise}
\begin{exercise}
Represent the standard generators of $H_2(S^2\times S^2)$ geometrically.
\end{exercise}
\begin{exercise}
Compute their mutual intersections.
\end{exercise}
\begin{exercise}
Compute their self-intersections using push-offs.
\end{exercise}
\begin{exercise}
Relate self-intersection to the normal Euler number.
\end{exercise}
\begin{exercise}
Explain why torsion is removed from the lattice form.
\end{exercise}
\begin{exercise}
Compare rank, parity, definiteness and signature as invariants.
\end{exercise}
\begin{exercise}
Write the intersection matrix for an orthogonal sum $H\oplus[1]$.
\end{exercise}
\begin{exercise}
Explain how VIII/34 prepares the graph--diagonal intersection picture in Lefschetz theory.
\end{exercise}

\section{Hints}
\begin{hint}
Use $H^2(S^4)=0$.
\end{hint}
\begin{hint}
Use the generator $h$ with $h^2[M]=1$.
\end{hint}
\begin{hint}
Reverse the fundamental class.
\end{hint}
\begin{hint}
Use the two factor classes.
\end{hint}
\begin{hint}
Find eigenvectors $(1,1)$ and $(1,-1)$.
\end{hint}
\begin{hint}
Count positive minus negative eigenvalues.
\end{hint}
\begin{hint}
Evaluate $Q((x,y),(x,y))$.
\end{hint}
\begin{hint}
Test the generator.
\end{hint}
\begin{hint}
Use orthogonal direct sum under connected sum.
\end{hint}
\begin{hint}
Combine $[1]$ and $[-1]$.
\end{hint}
\begin{hint}
Read signs from the diagonal.
\end{hint}
\begin{hint}
Replace $[M]$ by $-[M]$.
\end{hint}
\begin{hint}
Write the new basis matrix $P$.
\end{hint}
\begin{hint}
Use $\det(P)=\pm1$.
\end{hint}
\begin{hint}
Degree two has even parity.
\end{hint}
\begin{hint}
Degree one has odd parity.
\end{hint}
\begin{hint}
Use the two factor spheres.
\end{hint}
\begin{hint}
Perturb to transverse representatives.
\end{hint}
\begin{hint}
Push a factor sphere in its normal direction.
\end{hint}
\begin{hint}
Use the normal-bundle Euler class.
\end{hint}
\begin{hint}
Pair torsion with an integer-valued functional.
\end{hint}
\begin{hint}
Keep the invariants logically distinct.
\end{hint}
\begin{hint}
Write a block diagonal matrix.
\end{hint}
\begin{hint}
Think of fixed points as intersections of graph and diagonal.
\end{hint}
